\section*{Введение}
\addcontentsline{toc}{section}{Введение}

\textbf{Актуальность темы.} В условиях цифровой трансформации общества особую значимость приобретает автоматизация документооборота. Культурно-досуговые учреждения (дома культуры, театры, центры творчества) ежедневно обрабатывают значительные объемы документации, включая планы мероприятий, отчеты, договоры и внутренние служебные документы. Использование традиционных бумажных или частично автоматизированных методов ведения документации приводит к снижению эффективности работы, увеличению временных затрат на обработку информации и повышению риска ошибок. Разработка специализированных веб-приложений для автоматизации процессов документооборота позволяет повысить прозрачность управления, ускорить обработку данных и обеспечить централизованный доступ к информации. В качестве практической базы для анализа предметной области рассматривается деятельность МУ ЦКиД «Факел» г. Фрязино.

\textbf{Степень разработанности проблемы.} Вопросы автоматизации документооборота широко рассматриваются в научной и практической литературе. Существуют как универсальные системы электронного документооборота (например, «1С:Документооборот», «Directum», «DocsVision»), так и специализированные решения для отдельных отраслей. Однако большинство коммерческих систем характеризуется высокой стоимостью внедрения и избыточной функциональностью, не всегда соответствующей специфике культурно-досуговых учреждений. В то же время open-source решения требуют значительной доработки для адаптации под конкретные бизнес-процессы. Таким образом, задача разработки гибкого и адаптируемого веб-приложения для автоматизации документооборота в данной сфере остается актуальной.

\textbf{Объект исследования} — процессы организации и ведения документооборота в культурно-досуговых учреждениях (на примере МУ ЦКиД «Факел» г. Фрязино).

\textbf{Предмет исследования} — методы и программные средства автоматизации документооборота на основе веб-технологий.

\textbf{Целью работы} является проектирование и разработка веб-приложения для автоматизации документооборота в культурно-досуговых учреждениях на примере МУ ЦКиД «Факел» г. Фрязино.